\section{矩阵Lie群}
\label{sec:lie-matrix-groups}

置换群只有有限个元素，旋转却可以连续地变化：转角从$0$逐渐增大，旋转矩阵也随之连续变化。对这类群，可以问“单位元附近有哪些变化方向”，也可以要求群运算与这些方向的计算相容。本节先说明矩阵群怎样成为流形，并且乘法与求逆都是光滑映射；再用单位元处的切空间记录无穷小变化，即Lie代数；最后用矩阵指数把Lie代数中的方向变回群中的变换。

\begin{definition}[Lie群]
\label{def:lie-lie-group}
若群$G$同时是光滑流形，并且乘法$G\times G\to G$，$(g,h)\mapsto gh$与求逆$G\to G$，$g\mapsto g^{-1}$都是光滑映射，就称$G$为\term{Lie群}（Lie group）。
\end{definition}

定义要求群结构与光滑结构相容，而不只是“元素可以连续变化”。并非每个由连续变换组成的群都是有限维Lie群。例如以有理数倍$2\pi$为转角的旋转$\{\bm R(2\pi q):q\in\mathbb Q\}$是$\mathrm{SO}(2)$的子群，但它只有可数个元素，并且在单位圆上稠密：任意小的弧内都有它的元素。在继承的拓扑下，它的任一点附近既不像一段开区间，也不是孤立点，因而不是流形。圆周上全体同胚组成的群虽然也由连续变换组成，却不能用有限个实数参数局部描述，超出了本章讨论的范围。

本章只讨论矩阵组成的Lie群。称$\mathrm{GL}(n,\mathbb R)$的子群$G$为\term{矩阵Lie群}（matrix Lie group），若$G$在$\mathrm{GL}(n,\mathbb R)$中是闭集。可以证明，$\mathrm{GL}(n,\mathbb R)$的每个闭子群都是嵌入子流形，从而是Lie群（Cartan闭子群定理），这里不给出证明，可参考\scite{math-lee2013smooth}。下面用到的群都直接验证其流形结构，不依赖这个定理。上面的有理转角子群不是闭集，正好被这一定义排除。

\subsection{光滑结构}
\label{sec:lie-smooth-structure}

要把矩阵群当作流形，需要逐一说明它在矩阵空间中是什么样的子集、乘法与求逆为何光滑，以及各点的切空间怎样联系起来。下面依次处理一般线性群、正交群与左平移。

\paragraph{一般线性群}
把$n$阶实矩阵的$n^2$个元素排成一列，$\mathbb R^{n\times n}$就是欧氏空间。行列式是矩阵元素的多项式，因而连续，所以$\mathrm{GL}(n,\mathbb R)=\{\bm A:\det\bm A\neq0\}$是$\mathbb R^{n\times n}$中的开集，从而是$n^2$维流形，一张坐标图就覆盖全体。矩阵乘积的每个元素是两个因子元素的二次多项式，所以乘法光滑。由古典伴随公式$\bm A^{-1}=\operatorname{adj}(\bm A)/\det\bm A$，逆矩阵的元素是矩阵元素的有理函数，分母$\det\bm A$在$\mathrm{GL}(n,\mathbb R)$上不为零，所以求逆光滑。于是$\mathrm{GL}(n,\mathbb R)$是Lie群。

\paragraph{正交群}
正交群由约束$\bm A^\top\bm A=\bm I$给出。用第~\ref{chap:manifold-relations}章的正则水平集判断它是否是子流形时，约束映射的取值空间要选对。令$\mathrm{Sym}(n)$为$n$阶实对称矩阵组成的空间，维数为$n(n+1)/2$，定义
\[
\Phi:\mathrm{GL}(n,\mathbb R)\to\mathrm{Sym}(n),\qquad \Phi(\bm A)=\bm A^\top\bm A.
\]
$\bm A^\top\bm A$本来就对称，若把$\Phi$看成取值于$\mathbb R^{n\times n}$，上三角与下三角的方程会重复，微分不可能满秩。取值于$\mathrm{Sym}(n)$后，$\mathrm O(n)=\Phi^{-1}(\bm I)$。

\begin{theorem}[正交群的流形结构]
\label{thm:lie-orthogonal-group-manifold}
$\mathrm O(n)$是$\mathrm{GL}(n,\mathbb R)$中$n(n-1)/2$维的嵌入子流形，并且对$\bm A\in\mathrm O(n)$，
\begin{equation}
 T_{\bm A}\mathrm O(n)=\{\bm A\bm\Omega:\bm\Omega^\top=-\bm\Omega\}.
 \label{eq:lie-orthogonal-tangent}
\end{equation}
$\mathrm{SO}(n)$是$\mathrm O(n)$中的开子集，具有相同的维数与切空间。
\end{theorem}

\begin{proof}
沿方向$\bm H$求导，$\mathrm d\Phi_{\bm A}(\bm H)=\bm A^\top\bm H+\bm H^\top\bm A$。设$\bm A\in\mathrm O(n)$，任给$\bm S\in\mathrm{Sym}(n)$，取$\bm H=\frac12\bm A\bm S$，则
\[
\mathrm d\Phi_{\bm A}(\bm H)=\tfrac12\bm A^\top\bm A\bm S+\tfrac12\bm S^\top\bm A^\top\bm A=\tfrac12\bm S+\tfrac12\bm S=\bm S,
\]
所以$\mathrm d\Phi_{\bm A}$是满射，$\bm I$是正则值。由正则水平集定理，$\mathrm O(n)$是维数为$n^2-n(n+1)/2=n(n-1)/2$的嵌入子流形，切空间是$\ker\mathrm d\Phi_{\bm A}=\{\bm H:\bm A^\top\bm H\text{反对称}\}$。令$\bm\Omega=\bm A^\top\bm H$，由$\bm A\bm A^\top=\bm I$得$\bm H=\bm A\bm\Omega$，即式~\eqref{eq:lie-orthogonal-tangent}。最后，$\det$在$\mathrm O(n)$上连续且只取$\pm1$，所以$\mathrm{SO}(n)=\mathrm O(n)\cap\{\det>0\}$是$\mathrm O(n)$的开子集。
\end{proof}

$n=3$时维数为$3$，这正是三维朝向的自由度：九个矩阵元素满足六个独立的约束。乘法与求逆在$\mathrm O(n)$上光滑，是因为它们是$\mathrm{GL}(n,\mathbb R)$上光滑映射的限制，并且像仍落在$\mathrm O(n)$中；第~\ref{chap:manifold-relations}章说明，取值于嵌入子流形的光滑映射作为到该子流形的映射也是光滑的。$\mathrm{SL}(n,\mathbb R)=\det^{-1}(1)$可以同样处理：$\mathrm d\det_{\bm I}(\bm H)=\operatorname{tr}\bm H$不恒为零，在一般点上也满秩，所以它是$n^2-1$维的Lie群。

\paragraph{左平移}
式~\eqref{eq:lie-orthogonal-tangent}说明，$\bm A$处的切空间就是把单位元处的切空间整体左乘$\bm A$。这是群结构带来的一般现象。对Lie群$G$中的$g$，称$L_g:G\to G$，$L_g(h)=gh$为\term{左平移}（left translation）。它光滑，并且有光滑逆$L_{g^{-1}}$，所以是微分同胚，它的微分$\mathrm d(L_g)_e:T_eG\to T_gG$是线性同构。对矩阵群，$L_{\bm A}(\bm B)=\bm A\bm B$是线性映射，微分就是$\bm H\mapsto\bm A\bm H$。

因此只要知道单位元处的切空间，就能用群乘法得到任何一点处的切空间。用第~\ref{chap:manifold-basics}章建立的切空间语言，这把一个随点变化的对象化为一个固定的线性空间，下一小节就从这个线性空间入手。

\subsection{Lie代数}
\label{sec:lie-lie-algebra}

设$G$是矩阵Lie群。经过单位元的光滑曲线$\gamma:(-\varepsilon,\varepsilon)\to G$，$\gamma(0)=\bm I$，其速度$\gamma'(0)$是一个$n$阶矩阵，表示在单位元附近变换最先改变的方向。把这些速度全部收集起来，得到单位元处的切空间
\[
\mathfrak g=T_{\bm I}G=\{\gamma'(0):\gamma\text{为}G\text{中光滑曲线},\ \gamma(0)=\bm I\}.
\]
它是$\mathbb R^{n\times n}$的线性子空间。由上一小节，$\mathfrak{gl}(n)=\mathbb R^{n\times n}$是$\mathrm{GL}(n,\mathbb R)$的切空间，$\mathfrak{so}(n)=\{\bm\Omega:\bm\Omega^\top=-\bm\Omega\}$是$\mathrm O(n)$与$\mathrm{SO}(n)$的切空间，$\mathfrak{sl}(n)=\{\bm H:\operatorname{tr}\bm H=0\}$是$\mathrm{SL}(n,\mathbb R)$的切空间。

线性组合在$\mathfrak g$中有明确含义：若$\gamma_1,\gamma_2$是两条经过$\bm I$的曲线，则$\gamma(t)=\gamma_1(a t)\gamma_2(b t)$也在$G$中，由乘积求导得$\gamma'(0)=a\gamma_1'(0)+b\gamma_2'(0)$。所以把两个无穷小变化按比例叠加，对应于把两个小变换接连施加，二者在一阶上相同。

\paragraph{交换子}
接连施加的次序在一阶上看不出来，在二阶上却会显现。旋转一般不交换，Lie代数需要另一个运算记录这种不交换。

\begin{theorem}[交换子封闭]
\label{thm:lie-bracket-closed}
设$G$是矩阵Lie群，$\mathfrak g=T_{\bm I}G$。则对任意$\bm A\in G$与$\bm Y\in\mathfrak g$，有$\bm A\bm Y\bm A^{-1}\in\mathfrak g$；对任意$\bm X,\bm Y\in\mathfrak g$，有
\begin{equation}
 [\bm X,\bm Y]=\bm X\bm Y-\bm Y\bm X\in\mathfrak g.
 \label{eq:lie-commutator}
\end{equation}
\end{theorem}

\begin{proof}
设$\delta$是$G$中经过$\bm I$、速度为$\bm Y$的曲线，则$s\mapsto\bm A\delta(s)\bm A^{-1}$仍是$G$中经过$\bm I$的曲线，速度为$\bm A\bm Y\bm A^{-1}$，第一项得证。再设$\gamma$是速度为$\bm X$的曲线，令$c(t)=\gamma(t)\bm Y\gamma(t)^{-1}$。由第一项，$c(t)$始终落在线性子空间$\mathfrak g$中，所以它的导数也在$\mathfrak g$中。对$\gamma\gamma^{-1}=\bm I$求导得$(\gamma^{-1})'(0)=-\bm X$，于是$c'(0)=\bm X\bm Y-\bm Y\bm X$。
\end{proof}

称$[\bm X,\bm Y]$为\term{交换子}（commutator）。它满足三条规则：对每个变元线性；反对称，$[\bm X,\bm Y]=-[\bm Y,\bm X]$；Jacobi恒等式$[\bm X,[\bm Y,\bm Z]]+[\bm Y,[\bm Z,\bm X]]+[\bm Z,[\bm X,\bm Y]]=\bm0$，后者把三项展开后十二个乘积两两相消即得。一般地，带有满足这三条规则的二元运算的向量空间称为\term{Lie代数}（Lie algebra），矩阵Lie群的切空间$\mathfrak g$连同交换子是它的Lie代数。

\paragraph{与向量场括号的关系}
第~\ref{chap:manifold-basics}章用两个方向导数的交换差定义了向量场的Lie括号，这里的交换子与它一致。在开集$\mathrm{GL}(n,\mathbb R)\subseteq\mathbb R^{n\times n}$上，给定$\bm X$，令$\widetilde{\bm X}(\bm A)=\bm A\bm X$，这是由左平移把$\bm X$搬到每一点得到的向量场。欧氏开集上向量场的括号是$[V,W](\bm A)=\mathrm DW(\bm A)[V(\bm A)]-\mathrm DV(\bm A)[W(\bm A)]$；由$\mathrm D\widetilde{\bm Y}(\bm A)[\bm H]=\bm H\bm Y$，得
\[
[\widetilde{\bm X},\widetilde{\bm Y}](\bm A)=\bm A\bm X\bm Y-\bm A\bm Y\bm X=\bm A[\bm X,\bm Y].
\]
所以交换子正是这些左平移场的括号在单位元处的值。

\paragraph{旋转的Lie代数}
对$\bm\omega=(\omega_1,\omega_2,\omega_3)^\top\in\mathbb R^3$，令
\begin{equation}
 [\bm\omega]_\times=\begin{bmatrix}0&-\omega_3&\omega_2\\\omega_3&0&-\omega_1\\-\omega_2&\omega_1&0\end{bmatrix},
 \qquad [\bm\omega]_\times\bm x=\bm\omega\times\bm x.
 \label{eq:lie-hat-map}
\end{equation}
每个三阶反对称矩阵都唯一地写成这种形式，所以$\bm\omega\mapsto[\bm\omega]_\times$是$\mathbb R^3$到$\mathfrak{so}(3)$的线性同构，三维旋转的无穷小变化由一个三维向量记录。交换子对应叉积：$[[\bm a]_\times,[\bm b]_\times]=[\bm a\times\bm b]_\times$。验证时作用于任意$\bm x$，由二重叉积公式，
\[
\bm a\times(\bm b\times\bm x)-\bm b\times(\bm a\times\bm x)
=\bm b(\bm a^\top\bm x)-\bm a(\bm b^\top\bm x)=(\bm a\times\bm b)\times\bm x.
\]
例如$[[\bm e_1]_\times,[\bm e_2]_\times]=[\bm e_3]_\times$：绕第一轴与绕第二轴的小转动交替进行，二阶上产生绕第三轴的转动。

\paragraph{Lie代数元素与群元素}
Lie代数中的元素是速度，不是变换。$\mathfrak{so}(3)$中的矩阵反对称，而$\mathrm{SO}(3)$中的矩阵正交，二者除了作为矩阵之外没有包含关系：零矩阵在$\mathfrak{so}(3)$中，却不在$\mathrm{SO}(3)$中。Lie代数可以做线性组合，群却不能：$\bm I+\bm I=2\bm I$不是正交矩阵，两个旋转矩阵的平均一般也不是旋转矩阵。要从Lie代数回到群，需要下一小节的矩阵指数。

\subsection{矩阵指数}
\label{sec:lie-matrix-exponential}

给定$\bm X\in\mathfrak g$，我们希望找到一族变换$\gamma(t)\in G$，它以$\bm X$为初速度，并且“转$s$再转$t$等于转$s+t$”，正如$\bm R(\theta)$对转角那样。数的情形下，满足$f(s+t)=f(s)f(t)$、$f'(0)=x$的可微函数是$e^{tx}$。把指数函数的幂级数搬到矩阵上，就得到所需的变换。

\begin{definition}[矩阵指数]
\label{def:lie-matrix-exponential}
对$\bm X\in\mathbb R^{n\times n}$，定义
\begin{equation}
 \exp(\bm X)=\sum_{k=0}^\infty\frac{\bm X^k}{k!}=\bm I+\bm X+\frac{\bm X^2}{2}+\cdots,
 \label{eq:lie-matrix-exponential}
\end{equation}
称为$\bm X$的\term{矩阵指数}（matrix exponential）。
\end{definition}

级数对任意方阵绝对收敛。取一种满足次乘性$\lVert\bm A\bm B\rVert\leq\lVert\bm A\rVert\lVert\bm B\rVert$的矩阵范数，例如谱范数（可参考\BookChapterRef{chap:normed-inner-product-spaces}{卷二第3章“赋范空间与内积空间”}），则$\lVert\bm X^k/k!\rVert\leq\lVert\bm X\rVert^k/k!$，右边是收敛的数项级数$e^{\lVert\bm X\rVert}$。同一估计说明，在$\lVert\bm X\rVert\leq r$上级数一致收敛。对角矩阵的各次幂仍是对角矩阵，所以它的指数就是逐个对角元取指数。

\begin{lemma}[矩阵指数的基本性质]
\label{lem:lie-exponential-properties}
对$\bm X,\bm Y\in\mathbb R^{n\times n}$与可逆矩阵$\bm S$：
\begin{enumerate}
 \item 若$\bm X\bm Y=\bm Y\bm X$，则$\exp(\bm X+\bm Y)=\exp(\bm X)\exp(\bm Y)$；
 \item $\exp(\bm X)$可逆，$\exp(\bm X)^{-1}=\exp(-\bm X)$；
 \item $\exp(\bm X^\top)=\exp(\bm X)^\top$，$\exp(\bm S\bm X\bm S^{-1})=\bm S\exp(\bm X)\bm S^{-1}$；
 \item $t\mapsto\exp(t\bm X)$光滑，且$\frac{\mathrm d}{\mathrm dt}\exp(t\bm X)=\bm X\exp(t\bm X)=\exp(t\bm X)\bm X$；
 \item $\det\exp(\bm X)=e^{\operatorname{tr}\bm X}$。
\end{enumerate}
\end{lemma}

\begin{proof}
（1）$\bm X,\bm Y$交换时二项式定理成立，$(\bm X+\bm Y)^m=\sum_{k+l=m}\frac{m!}{k!\,l!}\bm X^k\bm Y^l$。两个绝对收敛级数的Cauchy乘积可以任意重排，所以
$\exp(\bm X)\exp(\bm Y)=\sum_{m}\sum_{k+l=m}\frac{\bm X^k\bm Y^l}{k!\,l!}=\sum_m\frac{(\bm X+\bm Y)^m}{m!}$。
（2）$\bm X$与$-\bm X$交换，由（1）得$\exp(\bm X)\exp(-\bm X)=\exp(\bm0)=\bm I$。
（3）逐项取转置或共轭，再用$(\bm S\bm X\bm S^{-1})^k=\bm S\bm X^k\bm S^{-1}$与级数的连续性。
（4）$t\bm X$与$h\bm X$交换，由（1）得$\exp((t+h)\bm X)-\exp(t\bm X)=(\exp(h\bm X)-\bm I)\exp(t\bm X)$，而
$\lVert(\exp(h\bm X)-\bm I)/h-\bm X\rVert\leq\sum_{k\geq2}\lvert h\rvert^{k-1}\lVert\bm X\rVert^k/k!\to0$，于是导数为$\bm X\exp(t\bm X)$；$\bm X$与$\exp(t\bm X)$交换，也可写成$\exp(t\bm X)\bm X$。重复这一论证得到各阶导数。
（5）令$f(t)=\det\exp(t\bm X)$。由（1）与行列式的乘法性，$f(t+h)=f(t)\det\exp(h\bm X)$。行列式在$\bm I$处的微分是迹，即$\det(\bm I+\bm H)=1+\operatorname{tr}\bm H+o(\lVert\bm H\rVert)$（可参考\BookChapterRef{chap:matrix-differentiation}{卷二第18章“矩阵函数的微分”}），而$\exp(h\bm X)=\bm I+h\bm X+O(h^2)$，所以$f'(t)=f(t)\operatorname{tr}\bm X$。于是$(e^{-t\operatorname{tr}\bm X}f(t))'=0$，由$f(0)=1$得$f(t)=e^{t\operatorname{tr}\bm X}$，取$t=1$即可。
\end{proof}

交换条件不能去掉。取$\bm X=\begin{bmatrix}0&1\\0&0\end{bmatrix}$、$\bm Y=\begin{bmatrix}0&0\\1&0\end{bmatrix}$，二者的平方都是零矩阵，所以$\exp(\bm X)=\bm I+\bm X$、$\exp(\bm Y)=\bm I+\bm Y$，乘积为$\begin{bmatrix}2&1\\1&1\end{bmatrix}$。而$\bm X+\bm Y$的平方是$\bm I$，按奇偶项汇总得$\exp(\bm X+\bm Y)=\cosh1\,\bm I+\sinh1\,(\bm X+\bm Y)$，对角元是$\cosh1\approx1.543\neq2$。

\paragraph{单参数子群}
由性质（1），$\exp((s+t)\bm X)=\exp(s\bm X)\exp(t\bm X)$，所以$t\mapsto\exp(t\bm X)$是从加法群$\mathbb R$到$\mathrm{GL}(n,\mathbb R)$的光滑同态，初速度为$\bm X$。这样的同态称为\term{单参数子群}（one-parameter subgroup）。下面的结论说明，初速度完全决定了单参数子群。

\begin{theorem}[单参数子群由初速度决定]
\label{thm:lie-one-parameter-subgroup}
设$\gamma:\mathbb R\to\mathrm{GL}(n,\mathbb R)$可微，$\gamma(s+t)=\gamma(s)\gamma(t)$对所有$s,t$成立，且$\gamma'(0)=\bm X$。则$\gamma(t)=\exp(t\bm X)$。
\end{theorem}

\begin{proof}
同态性给出$\gamma(0)=\bm I$与$\gamma(t+h)=\gamma(h)\gamma(t)$，所以$\gamma'(t)=\lim_{h\to0}\frac{\gamma(h)-\bm I}{h}\gamma(t)=\bm X\gamma(t)$。令$\bm Z(t)=\exp(-t\bm X)\gamma(t)$，由乘积求导与性质（4），
$\bm Z'(t)=-\bm X\exp(-t\bm X)\gamma(t)+\exp(-t\bm X)\bm X\gamma(t)=\bm0$，这里用了$\bm X$与$\exp(-t\bm X)$交换。于是$\bm Z(t)=\bm Z(0)=\bm I$，即$\gamma(t)=\exp(t\bm X)$。
\end{proof}

证明同时说明，$\bm Y(t)=\exp(t\bm X)\bm Y_0$是线性方程$\bm Y'=\bm X\bm Y$、$\bm Y(0)=\bm Y_0$的唯一解。线性常微分方程组的一般解理论可参考\BookChapterRef{chap:linear-differential-systems}{卷六第2章“线性微分方程与系统”}，这里只用到上面直接验证的部分。

\paragraph{Lie代数方向留在群内}
对矩阵Lie群$G$，若$\bm X\in\mathfrak g$，则对所有$t$有$\exp(t\bm X)\in G$；反过来，若对所有$t$有$\exp(t\bm X)\in G$，求导即知$\bm X\in\mathfrak g$。前一方向对一般闭子群的证明需要更多工具，这里不给出，可参考上文所引文献。对本章用到的群，可以直接验证。若$\bm\Omega$反对称，由性质（2）（3），
\[
\exp(\bm\Omega)^\top\exp(\bm\Omega)=\exp(\bm\Omega^\top)\exp(\bm\Omega)=\exp(-\bm\Omega)\exp(\bm\Omega)=\bm I,
\]
所以$\exp(\bm\Omega)\in\mathrm O(n)$；又由性质（5），$\det\exp(\bm\Omega)=e^{\operatorname{tr}\bm\Omega}=e^0=1$，所以$\exp(\bm\Omega)\in\mathrm{SO}(n)$。同理，迹为零的矩阵的指数行列式为$1$，落在$\mathrm{SL}(n,\mathbb R)$中。

例如$\bm J=\begin{bmatrix}0&-1\\1&0\end{bmatrix}\in\mathfrak{so}(2)$满足$\bm J^2=-\bm I$。把级数按奇偶项汇总，
\[
\exp(\theta\bm J)=\Bigl(1-\frac{\theta^2}{2!}+\cdots\Bigr)\bm I+\Bigl(\theta-\frac{\theta^3}{3!}+\cdots\Bigr)\bm J=\cos\theta\,\bm I+\sin\theta\,\bm J=\bm R(\theta),
\]
恰是式~\eqref{eq:lie-planar-rotation}。所以单参数子群$\theta\mapsto\bm R(\theta)$就是由速度$\bm J$生成的。

\paragraph{局部参数化的范围}
指数级数给出$\exp(\bm H)=\bm I+\bm H+O(\lVert\bm H\rVert^2)$，所以$\exp$在$\bm0$处的微分是恒等映射。把$\exp$限制在$\mathfrak g$上，它取值于$G$，微分$\mathfrak g\to T_{\bm I}G=\mathfrak g$可逆。由反函数定理，存在$\bm0$在$\mathfrak g$中的邻域$U$与$\bm I$在$G$中的邻域$V$，使$\exp:U\to V$是微分同胚。于是单位元附近的变换可以用Lie代数中的小向量作坐标；配合左平移，$\bm X\mapsto\bm A\exp(\bm X)$给出$\bm A$附近的坐标。

这个结论只在局部成立。指数映射一般不是单射：$\exp(2\pi\bm J)=\bm R(2\pi)=\bm I=\exp(\bm0)$。它也不一定是满射：由性质（5），$\exp(\bm X)$的行列式恒为正，所以$\mathrm{diag}(-1,1)\in\mathrm{GL}(2,\mathbb R)$不是任何实矩阵的指数。对$\mathrm{SO}(3)$，第~\ref{sec:lie-rotations}节会具体给出指数可逆的范围。

\paragraph{群指数与黎曼指数}
第~\ref{chap:riemannian-geometry}章的黎曼指数映射沿测地线走到时间$1$，依赖所选的度量与联络；这里的矩阵指数只依赖群的乘法。二者不能因名称相同而等同。把$\mathrm{GL}(n,\mathbb R)$看成$\mathbb R^{n\times n}$中带欧氏度量的开集，从$\bm I$出发、初速度为$\bm J$的测地线是直线$\bm I+t\bm J$，而单参数子群$\exp(t\bm J)=\bm R(t)$是一个圆；直线在$\det(\bm I+t\bm J)=1+t^2$下始终可逆，形状却与圆完全不同。

对$\mathrm{SO}(n)$配上$\mathbb R^{n\times n}$的Frobenius内积诱导的度量，两种指数在单位元处恰好一致。曲线$\gamma(t)=\exp(t\bm\Omega)$的二阶导数是$\gamma(t)\bm\Omega^2$，其中$\bm\Omega^2$对称；而$\gamma(t)$处的切向量形如$\gamma(t)\bm\Psi$，$\bm\Psi$反对称，$\langle\gamma\bm\Omega^2,\gamma\bm\Psi\rangle_{\mathrm F}=\operatorname{tr}(\bm\Omega^2\bm\Psi)$是对称矩阵与反对称矩阵之积的迹，等于零。环境加速度处处垂直于切空间，对带诱导度量的子流形，这正是测地线的条件。这一吻合依赖所选的度量与群的特殊对称性，不是一般结论。
